\documentclass[11pt]{article}

\usepackage[margin=1in]{geometry}
\usepackage{amsmath}
\usepackage{amssymb}
\usepackage{array}
\usepackage{booktabs}
\usepackage{caption}
\usepackage{enumitem}
\usepackage{float}
\usepackage{graphicx}
\usepackage{hyperref}
\usepackage{siunitx}
\usepackage{subcaption}
\usepackage{xcolor}

\graphicspath{{../../outputs/reports/figures/}}
\input{../../outputs/reports/generated/report_values.tex}

\hypersetup{
  colorlinks=true,
  linkcolor=blue!55!black,
  urlcolor=blue!55!black
}

\setlength{\parskip}{0.45em}
\setlength{\parindent}{0pt}

\title{Grid Transformation Report\\[0.4em]\large Axis-First Affine + Thin-Plate Spline Alignment}
\author{}
\date{\today}

\begin{document}
\maketitle

\section{Introduction}
We use one target speaker and one source speaker. Each speaker is represented by anatomical landmark points, 
and the goal is to map source points into target space. The default pair is target speaker and source speaker.

\begin{figure}[H]
  \centering
  \includegraphics[width=\textwidth]{raw_speakers.png}
  \caption{Raw target and source speaker images before any overlay or transformation.}
\end{figure}

\section{Element and Landmark Notation}
The speakers are denoted by landmark points on the superior tract, spine, jaw side, and posterior wall. These points are then used to build the grid.

\subsection*{Contours used to build the grid}
\begin{table}[H]
  \centering
  \begin{tabular}{>{\ttfamily}p{0.28\textwidth} p{0.60\textwidth}}
    \toprule
    Contour & Role in the grid \\
    \midrule
    incisior-hard-palate & Gives the superior hard-palate path and points \texttt{I1}--\texttt{I5}. \\
    soft-palate-midline & Gives the posterior superior points \texttt{I6} and \texttt{I7}. \\
    mandible-incisior & Gives the jaw-side points \texttt{M1} and \texttt{L6}. \\
    pharynx & Gives the posterior point \texttt{P1}. \\
    c1-c6 & Gives the spine points \texttt{C1}--\texttt{C6}. \\
    tongue & Only visualized and transferred; not used as a transform anchor. \\
    \bottomrule
  \end{tabular}
  \caption{Main contour elements used by the grid-transformation pipeline.}
\end{table}

\subsection*{Landmark and grid notation}
\begin{table}[H]
  \centering
  \begin{tabular}{>{\ttfamily}p{0.18\textwidth} p{0.70\textwidth}}
    \toprule
    Notation & Meaning \\
    \midrule
    I1--I7 & Superior tract points. \\
    C1--C6 & Spine points. \\
    M1, L6 & Jaw-side points. \\
    P1 & Posterior pharynx point. \\
    H1--H6 & Horizontal grid lines built from the landmark set. \\
    V1--V9 & Vertical grid lines sampled across the horizontal lines. \\
    \bottomrule
  \end{tabular}
  \caption{Landmark notation used in the report and in the code.}
\end{table}


\begin{figure}[H]
  \centering
  \includegraphics[width=\textwidth]{landmark_notation.png}
  \caption{Left: the main anatomical contours used to build the grid. Right: the landmark notation and the grid-line labels \texttt{H1}--\texttt{H6} and \texttt{V1}--\texttt{V9}.}
\end{figure}

\section{Speaker Differences and Why Two Steps}
The source and target speakers differ in both global pose and local tract shape. For that reason, the method first aligns the main axes with an affine transform and then refines the remaining mismatch with TPS.

\begin{figure}[H]
  \centering
  \includegraphics[width=\textwidth]{raw_with_grids.png}
  \caption{The same two speakers after grid construction. The grids make the anatomical differences explicit: the superior path, posterior wall, and tongue do not coincide even after simple visual normalization.}
\end{figure}

\section{Step 1: Affine Alignment}
Step 1 fits a full 2D affine map to the main grid axes. In the current implementation, it uses \StepOneAnchorCount{} axis anchors:
\par\smallskip
\StepOneAnchorList.

The transform is
\begin{equation}
\mathbf{x}' = A\mathbf{x} + \mathbf{t},
\end{equation}
where \(A \in \mathbb{R}^{2 \times 2}\) models rotation, anisotropic scale, and shear, and \(\mathbf{t} \in \mathbb{R}^{2}\) models translation. The parameters are obtained by least squares:
\begin{equation}
\min_{A,\mathbf{t}} \sum_{i=1}^{N} \left\lVert A\mathbf{x}_i + \mathbf{t} - \mathbf{y}_i \right\rVert^2.
\end{equation}

\begin{figure}[H]
  \centering
  \includegraphics[width=\textwidth]{step1_affine.png}
  \caption{Step 1 affine alignment. Left: the source grid before alignment. Middle: affine-mapped anchor points against their target partners. Right: the source grid after the global affine fit.}
\end{figure}

\section{Step 2: Thin-Plate Spline Refinement}
Let \(\{\mathbf{x}_i\}_{i=1}^{N}\) be the post-affine source control points and \(\{\mathbf{y}_i\}_{i=1}^{N}\) be the target control points. In the current implementation, Step 2 uses \StepTwoControlCount{} controls:
\par\smallskip
\StepTwoControlList.

The TPS output is a non-linear map \(f=(f_x, f_y)\) such that \(f(\mathbf{x}_i)=\mathbf{y}_i\). Each output coordinate is modeled as:
\begin{align}
f_x(\mathbf{p}) &= a_1 + a_2 p_x + a_3 p_y + \sum_{i=1}^{N} w_i^{(x)} U\!\left(\lVert \mathbf{p} - \mathbf{x}_i \rVert \right), \\
f_y(\mathbf{p}) &= b_1 + b_2 p_x + b_3 p_y + \sum_{i=1}^{N} w_i^{(y)} U\!\left(\lVert \mathbf{p} - \mathbf{x}_i \rVert \right),
\end{align}
with the 2D thin-plate-spline kernel
\begin{equation}
U(r) = r^2 \log r.
\end{equation}

In matrix form, the TPS coefficients solve
\begin{equation}
\begin{bmatrix}
K & P \\
P^\top & 0
\end{bmatrix}
\begin{bmatrix}
\mathbf{w} \\
\mathbf{a}
\end{bmatrix}
=
\begin{bmatrix}
\mathbf{d} \\
\mathbf{0}
\end{bmatrix},
\end{equation}
where \(K_{ij} = U(\lVert \mathbf{x}_i - \mathbf{x}_j \rVert)\) and \(P\) contains the affine tail \([1, x_i, y_i]\).

The final output map is
\begin{equation}
T(\mathbf{p}) = f(A\mathbf{p} + \mathbf{t}).
\end{equation}

\begin{figure}[H]
  \centering
  \includegraphics[width=\textwidth]{step2_tps.png}
  \caption{Step 2 TPS. Left: control-point correspondences after Step 1. Right: the local bending effect from the affine result to the final source grid.}
\end{figure}

\begin{figure}[H]
  \centering
  \includegraphics[width=0.82\textwidth]{final_alignment.png}
  \caption{Final source-to-target alignment after the composed affine + TPS transform. The target tongue is shown in solid pink and the mapped source tongue in dashed pink.}
\end{figure}

\section{Applying the Final Transform to Tongue and Image}
\subsection*{Contours such as tongue}
Once the final transform \(T\) has been estimated, any open contour can be moved point-by-point. If \(C_{\text{src}}(s)\) is a source contour parameterized by arc position \(s\), the moved contour is
\begin{equation}
C_{\text{moved}}(s) = T\!\left(C_{\text{src}}(s)\right).
\end{equation}

In the current pipeline, the same point mapping is applied to the shared articulators used in the overlays: hard palate, soft-palate midline, soft palate, tongue, mandible-incisior, and pharynx. After mapping, the repository applies a short endpoint-preserving smoothing pass for visualization only:
\begin{equation}
\widetilde{C}_{\text{moved}}(s) = S\!\left(C_{\text{moved}}(s)\right).
\end{equation}
This smoothing does \emph{not} change the affine fit or the TPS coefficients; it only reduces jagged visual artifacts on the displayed contours. For the current pair, the tongue RMS after transfer is \TongueRMS{} px.

\begin{figure}[H]
  \centering
  \includegraphics[width=\textwidth]{tongue_transfer.png}
  \caption{Applying the final transform to articulator contours. The tongue is highlighted explicitly and transferred with the same point transform used for the grid.}
\end{figure}

% \subsection*{Full image warp}
% The full source speaker image is warped into target space with inverse mapping. For each target-space pixel \((x_t, y_t)\), the code evaluates the inverse transform and samples the source image at the corresponding source-space position:
% \begin{equation}
% I_{\text{warp}}(x_t, y_t) = I_{\text{src}}\!\left(T^{-1}(x_t, y_t)\right).
% \end{equation}
% The current implementation performs this sampling with \texttt{scipy.ndimage.map\_coordinates} using bilinear interpolation (\texttt{order=1}). This is the same geometric transform as the contour transfer, but applied densely to every image pixel instead of only to contour vertices.

\begin{figure}[H]
  \centering
  \includegraphics[width=\textwidth]{full_image_warp.png}
  \caption{Applying the final transform to the full source image. The warped source frame is shown in target space together with articulator overlays for comparison.}
\end{figure}

% \section{Conclusion}
% The current grid-transformation pipeline is deliberately modular and interpretable:
% \begin{itemize}[leftmargin=1.4em]
%   \item Step 1 affine solves the global anatomical pose mismatch.
%   \item Step 2 TPS refines local residual mismatch without discarding the axis-first alignment.
%   \item The same final transform can then be reused for grids, articulator contours, tongue transfer, and dense image warping.
% \end{itemize}
%
% The main limitation is that the tongue is still a transferred \emph{payload}, not an alignment anchor. This keeps the grid fit anatomically grounded in the palate, jaw, pharynx, and spine axes, but it also means residual tongue mismatch can remain when tongue shape changes independently of those anchors.

\end{document}
